\section{The Map and Capacity (The Dirac Propagator)}
\label{sec:map_capacity}

\textbf{The Architectural Requirement:}
\begin{itemize}
    \item \textbf{The Persistence Problem:} Information must propagate across the causal manifold without amplitude loss (Map), but it cannot travel strictly at the speed of light without completely delocalizing and losing its localized structural memory (Capacity).
    \item \textbf{The Physics Result:} The Dirac Lagrangian ($\mathcal{L} = \bar{\psi}i\gamma^\mu D_\mu\psi - m\bar{\psi}\psi$). The kinetic term is the unique geometric penalty enforcing unitary chiral propagation, and the mass term provides the scalar impedance necessary to anchor the state into a persistent local frame.
\end{itemize}

\textbf{Terminology Mapping:} Throughout this text, we standardize the terminology bridging geometry and physics. A localized, persistent boundary structure on the substrate (characterized by the topological invariant $\chi=2$) is a \textbf{Topological Defect}. When this defect occupies the chiral node channels ($\nu=16$), its physical manifestation is a \textbf{Fermion}. 

We first derive the dynamics of the \textbf{Map} pillar by defining the evolution of a fermion state $\psi(x,t)$ on the lattice. Time is treated not as a continuous dimension, but as a discrete clock cycle defined by the fundamental temporal step $\tau$.

\subsection{The Unitary Flux Constraint}
\label{sec:unitary_flux_constraint}

The Map pillar requires that information propagate locally while strictly conserving probability amplitude (Unitarity). The uniquely appropriate discretization scheme to enforce the three core architectural requirements derived for the vacuum: locality (nearest-neighbor transitions), unitarity (probability conservation at each discrete step), and finiteness (a strictly bounded state space per node), is structurally equivalent to the \textbf{Quantum Lattice Gas Automaton (QLGA)} formalism \cite{meyer_quantum_1996}\footnote{The QLGA formalism provides a rigorous mathematical framework for unitary, local, discrete-time evolution on a lattice. Our discrete update rule is structurally equivalent to a QLGA on the $D_4$ geometry. Crucially, it is \emph{not} a Wilson-fermion discretization. Standard Wilson fermions suppress doublers by introducing an ad hoc dimension-5 operator ($a \nabla^2$), which strictly violates our Equipartition of Entropic Cost (\cref{subsec:mapping_rules}) by exceeding the node capacity limit ($\nu=16$). Instead, the $E_8$-Persistence theory lattice evades fermion doubling purely geometrically: the orthogonal projection strictly isolates the mirror chiralities into the static internal sector, requiring no continuum-limit cancellations.}, generalizing the Feynman Checkerboard model\cite{feynman_quantum_1965} to the $D_4$ geometry.

\textbf{The Discrete Update.} The fermion state $\psi$ (a 16-component chiral spinor occupying the node capacity $\nu=16$) updates via a unitary operator combining two actions:

\begin{enumerate}
    \item \textit{Shift:} Information propagates to nearest neighbors with geometric admittance (hopping amplitude) $w$ at speed $c = \ell/\tau = 1$, mediated by internal spatial transition matrices $\Gamma^k$. The unit vector $e_k$ defines the $k$-th spatial axis. The symmetric difference across these neighbors encodes the local directional gradient, which is scaled by $\Gamma^k$ and $w$, and phase-rotated by $-i$ to preserve unitarity.

    \item \textit{Mass:} A local rotation mixing chiral components with topological impedance (mass) $m$, phase-rotated by $-i$, and governed by the temporal projection matrix $\gamma^0$ (the operator defining the temporal update of the local spinor state). (As we will rigorously prove in \cref{sec:capacity_mechanism}, this impedance $m$ cannot be a rigid bare constant due to orthogonal sub-lattices, but must dynamically emerge from the Governor field. We will also formally identify this geometric projection matrix $\gamma^0$ as identically the $\gamma^0$ of the continuum Dirac equation).
\end{enumerate}

The discrete evolution operator acts as:
\begin{equation}
\begin{aligned}
    \psi(x,t+\tau) &= \psi(x,t) \\
    &\quad - iw \sum_{k=1}^3 \Gamma^k \left[\psi(x+e_k\ell,t) - \psi(x-e_k\ell,t)\right] \\
    &\quad - i\gamma^0 m\tau\psi(x,t)
\end{aligned}
\label{eq:unitary_flux}
\end{equation}

To recover the macroscopic field equation, we map this discrete update rule to its continuous Effective Field Theory (EFT) approximation. Subtracting $\psi(x,t)$ and dividing by $\tau$ expands the central difference quotients to first order:
\begin{equation}
\begin{aligned}
    \frac{\psi(x,t+\tau) - \psi(x,t)}{\tau} = \\
    -iw \sum_{k=1}^3 \Gamma^k \left( \frac{\psi(x+e_k\ell,t) - \psi(x-e_k\ell,t)}{2\ell} \right) &\frac{2\ell}{\tau} - i\gamma^0 m\psi(x,t)
\end{aligned}
\label{eq:diff_form}
\end{equation}

\textbf{The Ontological Reversal:} Crucially, the lattice spacing $\ell$ and time step $\tau$ do not physically vanish. They represent the finite, immutable Planck-scale resolution of the substrate. Unlike standard QFT where the continuous manifold is assumed fundamental and the discrete lattice is merely a mathematical regularization trick, $E_8$-Persistence theory reverses this hierarchy: the discrete substrate is the exact physical reality, and the continuous differential manifold is strictly a macroscopic approximation.

\textbf{The Geometric Constraint.} Standard lattice formulations treat the hopping amplitude as a tunable parameter. Here, the kinematic bandwidth limit $c \equiv \ell/\tau = 1$ strictly enforces $w = 1/2$. To see this, note that the central difference quotient spans $2\ell$ (from $-\ell$ to $+\ell$), so mapping this to the continuous spatial derivative $\partial_k\psi$ requires the discrete prefactor $-iw(2\ell/\tau) = -i(2wc)$ to identically match the continuum Dirac velocity coefficient ($-i$). Thus $2w = 1$, fixing $w = 1/2$ with zero free parameters.

Identifying $\partial_t$ and $\partial_k$ in the continuum limit, the discrete evolution maps uniquely to the continuous Hamiltonian form:
\begin{equation}
    i \partial_t \psi = -i \sum_{k=1}^{3} \Gamma^k \partial_k \psi + \gamma^0 m \psi
\end{equation}
The precise geometric nature and algebraic properties of the transition matrices ($\Gamma^k, \gamma^0$) are uniquely determined by the causal geometry of the lattice, as derived below.

\subsection{The Clifford Identification}

\textbf{Architectural Requirement:} The \textbf{Map} pillar requires lossless propagation of a complex signal with precise directional resolution across the discrete substrate.

\textbf{Lattice Constraint:} The transition matrices ($\Gamma_\mu$) mediating nearest-neighbor updates must simultaneously satisfy two architectural requirements: Unitarity (probability conservation at each discrete step) and metric compatibility (reproducing the Lorentzian causal structure derived from the $E_8 \to D_4 \oplus D_4$ projection \cite{meyer_informational_2026}). 

To ensure isotropic propagation, the cross-terms of the discrete spatial transitions must identically vanish. The unique algebraic relation satisfying both the Unitarity constraint and the metric signature is the anti-commutation relation:
\begin{equation}
    \{\Gamma^\mu, \Gamma^\nu\} \equiv \Gamma^\mu\Gamma^\nu + \Gamma^\nu\Gamma^\mu = 2\eta_{\mu\nu}\mathbb{I}
\label{eq:anti_commute}
\end{equation}
where $\eta_{\mu\nu} = \operatorname{diag}(-1,+1,+1,+1)$ is the emergent Lorentzian metric, geometrically enforcing $(\Gamma^0)^2 = -\mathbb{I}$ (temporal inversion) and $(\Gamma^k)^2 = +\mathbb{I}$ (spatial reflection) for spatial indices $k$.

\textbf{Recognition:} This anti-commutation relation is the strict geometric requirement for preserving causal structure on the lattice, not a postulated quantum mechanical spin algebra. Because \cref{eq:anti_commute} is identically the defining algebra of the Dirac matrices, we explicitly identify the internal lattice transition matrices with the physical gamma matrices ($\Gamma^\mu \equiv \gamma^\mu$). Furthermore, this unique algebraic relation preserves the Unitarity of the discrete shift operator while strictly suppressing lattice doublers through the geometric isolation of chiralities established in the substrate derivation\cite{meyer_informational_2026}.

To recover the covariant Dirac equation, we multiply the continuum equation by $\gamma^0$ (utilizing the Lorentzian metric property $(\gamma^0)^2 = -1$), which perfectly recovers the covariant form:
\begin{equation}
    (i\gamma^\mu \partial_\mu - m)\psi = 0
    \label{eq:covariant_dirac}
\end{equation}

\subsection{The Chiral Filter: Geometric Origin of Parity Violation}
Standard QFT inserts the chiral projection $P_L = \frac{1-\gamma^5}{2}$ by hand to match the observed maximal parity violation of the weak interaction. Here, this asymmetry emerges structurally from the substrate geometry.

The \textbf{Map} pillar requires information to propagate on the causal (Lorentzian) manifold. The geometric projection $E_8 \to D_4 \oplus D_4$ splits the 16-component spinor into two orthogonal sectors. Because the Lorentzian signature assigns the active temporal generator ($\partial_t$) exclusively to the spacetime sector, the corresponding internal symmetry sector is strictly Euclidean. 

Consequently, the Dirac \cref{{eq:covariant_dirac}} automatically decomposes. The left-handed components evolve actively under the temporal generator, while the right-handed components remain static with respect to the internal symmetry's active temporal flow. 

Thus, the bare internal gauge generators necessarily couple only to the left-handed projection $P_L\psi$. Maximum parity violation is enforced by the geometric absence of a temporal update path for the right-handed sector within the internal symmetry space, rendering right-handed fermions strictly as singlets under this bare interaction (the structural origin of the Weak interaction's $SU(2)_L$ chirality).

\subsection{The Capacity Mechanism: Mass as Chiral Impedance}
\label{sec:capacity_mechanism}

If the vacuum only possessed the Map pillar, the resulting massless Dirac equation ($\mathcal{L}_{Map} \propto \bar{\psi} i\gamma^\mu D_\mu \psi$) would describe pure ballistic transport at $v=c$. While this perfectly satisfies information propagation, it violates the \textbf{Capacity} pillar. A signal moving at the speed of light possesses no rest frame; it is entirely delocalized along its light cone and thus cannot store local configuration memory.

\textbf{The Geometric Clash:} To satisfy Capacity, the state must arrest its spatial translation, establishing a local address. In the spinor formalism, this requires a resonant loop: continuously converting an outgoing Left-Chiral flux into an incoming Right-Chiral flux (a literal geometric realization of Schr\"{o}dinger's \textit{Zitterbewegung}).

However, as derived in the previous section, the $E_8 \to D_4 \oplus D_4$ projection strictly isolates the active spacetime kinematics (Left) from the static internal symmetry (Right) into orthogonal sub-lattices. A bare mass term ($-m\bar{\psi}\psi = -m(\bar{\psi}_L\psi_R + \bar{\psi}_R\psi_L)$) represents a direct, rigid cross-term between orthogonal spaces. This is geometrically forbidden. This structural prohibition recovers the chiral gauge invariance of the Standard Model not as an abstract symmetry, but as a rigid geometric absolute.

\textbf{The Scalar Bridge and the Dimensional Budget:} How does the system bridge these orthogonal sectors to achieve Capacity? We return to the Dimensional Budget established in \cref{subsec:restmass}. A fermion node intrinsically consumes 3 units of the dimensional budget. To persist across continuous temporal updates without spatial translation (kinetic flux $f \to 0$), it must couple to exactly one dimension-1 structural parameter to exactly saturate the Equipartition of Entropic Cost ($[K] + [f] = 4$).

Because a static bare mass ($m$) is forbidden by the orthogonal geometry, this 1-dimensional bridge cannot be a rigid constant. It mathematically must be a dynamic, phase-carrying scalar field capable of mediating momentum and quantum numbers across the orthogonal boundary. Rather than an ad-hoc invention, this scalar field is the physical manifestation of the Governor pillar (the constraint mechanism preventing unbounded divergence, \cref{sec:governor}).

\textbf{The Yukawa Operator:} To arrest information flow and satisfy Capacity, the geometry uniquely mandates that the fermion couples to this Governor field ($\phi$). The bare $\mathcal{L}_{CAP}$ term must therefore take the form of a Yukawa Coupling:
\begin{equation}
    \mathcal{L}_{CAP} = -y \left( \bar{\psi}_L \phi \psi_R + \bar{\psi}_R \phi^\dagger \psi_L \right) \label{eq:yukawa}
\end{equation}
where $y$ is the geometric admittance of the chiral bridge, fixed by the substrate coupling efficiency ($g=1/2$) and embedding geometry.

In Quiescent Equilibrium, where the Governor field minimizes its own entropic action by settling to a constant non-zero background expectation value ($\langle\phi\rangle$), this dynamic bridging operator reduces exactly to the effective mass impedance term presented in the Entropic Balance Sheet: $\mathcal{L}_{CAP} \to -m\bar{\psi}\psi$, where $m = y\langle\phi\rangle$.

\textbf{Conclusion: The Physical Nature of Mass.} The Dirac equation is not an axiom; it is the hydrodynamic limit of a discrete, unitary cellular automaton evolving on a spinor lattice. Fermion mass ($m = y\langle\phi\rangle$) is the geometric impedance required to arrest a state from dissipating at the speed of light, locking a topological defect into a persistent local capacity.